\section{A verified family without central symmetry}
\label{sec:v36-nonelliptic-family}
The local and path-bridge principles also apply beyond quadratic
scatterers. We give an explicit fixed-degree family on the same
triangular lattice. The family need not have central symmetry, so it
uses the rotation-free finite-cover argument essentially.

Fix an integer $D\ge3$. For parameters in the compact set
\begin{equation}\label{eq:v36-support-parameters}
 \frac{91}{200}\le r\le\frac{93}{200},\qquad
 \sum_{j=2}^D(j^2-1)\sqrt{a_j^2+b_j^2}\le\frac1{200},
\end{equation}
define a support function and its normal parametrization by
\begin{align}
 h(\theta)&=r+\sum_{j=2}^D(a_j\cos j\theta+b_j\sin j\theta),\notag\\
 x(\theta)&=h(\theta)n_\theta+h'(\theta)t_\theta.
                                                \label{eq:v36-support-curve}
\end{align}
The index $j=1$ is omitted to fix translations. Arbitrary orientation
is already permitted by the paired sine and cosine coefficients.
The positive quantity $h+h''$ is the radius of curvature, not the
radial distance from the chosen center.

\begin{proposition}[Nonelliptic action families]
\label{prop:v36-support-family}
Place identical translates of the convex body in
\eqref{eq:v36-support-curve} at the triangular lattice points.
For each fixed $D$, the full compact family
\eqref{eq:v36-support-parameters} satisfies \textup{(H1)--(H7)}.
Consequently its stationary microscopic local law, regular selected
law, exact conditional physical path bridge and positive-band
posterior comparison all hold uniformly in its parameters.
Its perimeter and enclosed area are
\begin{equation}\label{eq:v36-support-area}
 L=2\pi r,\qquad
 A=\pi r^2-\frac\pi2\sum_{j=2}^D(j^2-1)(a_j^2+b_j^2),
 \qquad \bar\tau=\frac{\pi(\sqrt3/2-A)}{L}.
\end{equation}
There are members of this family which are not centrally symmetric
and hence are not ellipses.
\end{proposition}
\begin{proof}
Cauchy--Schwarz applied to each sine--cosine pair and the coefficient
condition give
\[
 \frac9{20}\le h+h''\le\frac{47}{100},\qquad
 |h-r|\le\frac1{600}.
\]
Thus the curve is a smooth strictly convex embedded oval, and its
body contains the centered disk of radius $9/20$ and is contained in
the disk of radius $47/100$. Indeed disk containment is equivalently
comparison of support functions. Distinct obstacles have separation
at least $3/50$; the contained smallest disks supply the uniform finite
horizon. Curvatures have common positive upper and lower bounds, and
fixed degree gives common bounds for every fixed boundary derivative.
Since $x'=(h+h'')t_\theta$, normalized arclength and its inverse have
common smooth bounds. The unweighted local-space theorems and
Proposition~\ref{prop:v36-cover-mixing} establish (H2)--(H3), and
the general flight action establishes (H4).

We verify the two regularity issues not implied by curvature bounds
alone. In rational half-angle charts, $x(\theta)$ and $n_\theta$ are
rational functions with uniformly nonvanishing chart denominators;
$r,a_j,b_j$ are parameters. The backwards velocity
$\cos\varphi\,n_\theta-\sin\varphi\,t_\theta$ is treated in the
same charts. A candidate first intersection with a preceding obstacle
can be specified by the existence of its boundary angle and a
positive flight length, together with nonintersection at smaller
lengths. Quantifier elimination gives semialgebraic sets of bounded
complexity, depending on fixed $D$ but not on a parameter value.
Finite horizon leaves only a fixed finite candidate list.
Uniform cell decomposition therefore yields boundedly many regular
first-hit regions and monotone boundary arcs. Ties of two regular
first hits cannot occur for separated obstacles. The genuine
boundaries are backwards collision singularities, in the unstable
cone, and have uniform transverse slope separation from the stable
curves. The same finite chart refinement as in the ellipse proof
therefore gives bounded stable intersections and the
$O(\epsilon)$ boundary-neighborhood length estimate in every
homogeneity strip. This decomposition is done before normalizing
arclength; a smooth monotone change then preserves these properties.

The flight length is uniformly piecewise $C^{1/2}$, including in the
parameters. To see the exponent explicitly, take the signed-distance
function to a candidate obstacle in a common tubular neighborhood.
Along a tangent straight ray its first derivative vanishes and its
second derivative is the strictly positive curvature. These second
derivatives are bounded away from zero, uniformly on the compact
family. The parameter-dependent one-variable Morse normal form
therefore expresses the two intersection branches, near tangency, as
a smooth change of variable of $\pm\sqrt{b}$, where $b\ge0$ is a
smooth parameter-and-state function. Uniform local derivatives and
$|\sqrt b-\sqrt c|\le\sqrt{|b-c|}$ give the claimed bound.
A finite cover of the compact tangency set supplies common constants.
Away from these neighborhoods the entering intersection is simple
and the implicit-function theorem gives uniform smooth bounds.
Only the entering branch with least positive flight length is used;
the current obstacle's zero root is discarded. This proves (H5)
and the matched-root part of (H6). The remaining one-step parameter
comparison is the collision-space comparison in (H2), with these
root bounds inserted, as in Section~\ref{sec:action-norm-details}.

For (H7), a bounded straight flight intersects a convex polygonal
cell in an interval and visits only uniformly finitely many cells.
Its interval endpoints vary uniformly off the singular, tangent and
vertex-crossing sets. In the preceding rational charts these sets
are semialgebraic of fixed complexity. They are finite unions of
null arcs and points, including possible edge-coincident flights:
the latter have a prescribed velocity direction, a codimension-one
condition in $(\theta,\varphi)$. Closed neighborhoods at continuity
radii have null boundary and arbitrarily small limiting measure.
On their compact complements the regular first roots and cell
intersections persist along any convergent parameter sequence.
This verifies the precise closed-exceptional-neighborhood condition,
not just pointwise continuity of intervals.

Finally $L=\int_0^{2\pi}(h+h'')=2\pi r$ and
$A=\tfrac12\int_0^{2\pi}h(h+h'')$. Fourier orthogonality gives
\eqref{eq:v36-support-area}; the mean roof is the same exact suspension
identity as before. For example, $h(\theta)=23/50+\epsilon\cos3\theta$
with $0<|\epsilon|\le1/1600$ satisfies the coefficient constraint.
A translation changes only the first harmonic of a support function;
a nonzero third harmonic cannot be removed that way. Hence this body
is not centrally symmetric. The conclusions follow from
Theorem~\ref{thm:v35-action-family-LLT},
Theorem~\ref{thm:v36-microscopic-path-bridge}, and
Corollaries~\ref{cor:v35-family-posterior} and
\ref{cor:v36-band-posterior-bridge}.
\end{proof}

No uniform assertion is made as $D\to\infty$. The degree is fixed
before the finite-complexity constants are chosen. This family is an
additional application on the original lattice, not a deformation of
the four-coordinate return observable or its section.
